% ============================================================
% DRAFT — Appendix B: Graduate Attribute Outcome Compliance.
%
% Not something the framework spells out as mandatory text, but every
% past report I could check (Wait 2020, Baas 2020, Theron 2025) has
% exactly this appendix, and it's the easiest way to hand the examiner
% a quick-reference tickbox. Required GAs, confirmed from your own
% signed PE448 2026 framework PDF (Project E 448_2026_Framework_Eng_
% signed.pdf): GA 1, 2, 3, 4, 5, 6, 8, 9. GA 7 is NOT on your module's
% list -- don't add it, every reference report I checked also skips it.
%
% IMPORTANT things to check before you use this:
%   1. Section references below are written as plain chapter/section
%      NAMES, not \ref{} labels, because I don't have your final label
%      names for Design/Implementation (and you're still restructuring
%      those -- cutting Table 3.5, maybe merging the Hardware chapter).
%      Swap each one for a real \ref{} once your structure is locked,
%      or just leave them as prose if that's easier.
%   2. GA 6 explicitly requires BOTH the written report AND a
%      satisfactory oral exam, per the framework text ("GA 6: The
%      project report (written communication) AND a satisfactory oral
%      (oral communication)") -- that's different from GA 1-5/8/9,
%      which the report alone can satisfy. Said explicitly below.
%   3. GA 8 (Individual Work): I have NOT written a specific claim
%      about your own use of AI tools in drafting this report -- that's
%      a factual statement about your own process, not mine to write
%      for you. I left a placeholder. Worth knowing: Theron's 2025
%      report (one you just had me read) states this plainly in its
%      own GA8 entry: "The use of AI tools was used to refine language
%      in the report, but the content itself was done by the student."
%      If that's roughly true for you too, say so in your own words --
%      an honest, specific disclosure here is safer than silence,
%      given how much drafting work has gone through this chat.
%   4. GA 9: your own status notes (16 Sept) flagged this one as "still
%      needs a project-specific draft" -- meaning you hadn't started it
%      yet as of then. I've drafted a first pass below from what I know
%      of the project, but check it's actually true to your experience.
%   5. Every example report I read is specific -- it names actual
%      libraries, algorithms and numbers, not generic claims. I've
%      tried to match that (Ollama, Gemma 4 E4B, ROS2/rosbridge, the
%      actual run counts) rather than write something generic that
%      could apply to any project. Keep that specificity when you edit.
%
% This is a draft for review. Nothing here has been written into your
% actual chapter/appendix files.
% ============================================================

\chapter{Graduate Attribute Outcome Compliance}
\label{app:ga}

All required ECSA exit-level outcomes were achieved while executing this
project and writing this report. This section describes how each one was
achieved, with references to where the evidence appears in the report.

\section*{GA 1. Problem Solving}

The core problem -- turning an ambiguous spoken command into a safe,
executable plan for a resource-constrained rover -- is genuinely
under-specified: commands range from a single literal action to a vague,
multi-step request with no fixed target (Table~1.1, command levels L0--L5).
I had to define my own taxonomy of command complexity and a bounded
capability vocabulary to make the problem tractable at all (Chapter~1,
Chapter~3 System Design). Solving it meant working within real constraints
that forced creative compromises rather than textbook solutions: no onboard
compute budget for a large model, an unreliable vision-based distance
estimate that had to be abandoned mid-project in favour of a fixed hop
length once testing showed it wasn't usable (Chapter~6, ``Distance from
vision''), and a rover with no wheel encoders, which shaped the
stop-look-hop and ultrasonic-gating approach (Chapter~3, Chapter~5).

\section*{GA 2. Application of Scientific and Engineering Knowledge}

Mathematics and statistics were used to define the correctness metrics used
throughout Chapter~6 (\texttt{visits\_match}, \texttt{first\_divergence},
per-depth and per-context breakdowns over 270 virtual runs). Sensor and
control reasoning underpins the ultrasonic-stop threshold and the rejected
vision-distance approach (Chapter~6, ``Executability on the Physical
Chassis''). Computer science and machine learning fundamentals --
prompt construction, local LLM inference behaviour, context-window
capacity versus generation-length limits -- underpin the planner design
and are reasoned about explicitly in Chapter~6 rather than assumed.

\section*{GA 3. Engineering Design}

The overall system architecture is a deliberate, justified design choice:
compute-heavy planning and vision run on a laptop/desktop (``the brain''),
while the TurboPi is treated as instrumentation only (Chapter~1 Background,
Chapter~3, Figure~3.1 ``The compute split''). Alternatives were considered
and explicitly rejected with reasons, not just assumed away: an onboard
Jetson was rejected on cost/compute/generality grounds (Chapter~1); a
variable hop length driven by the vision model's distance estimate was
designed, tested, and then rejected once testing showed neither vision
model gave a usable absolute distance (Chapter~6). Component-level design
work includes the capability vocabulary, the multi-turn clarification state
machine (Chapter~3, Figure~3.3), and the guidance/recovery loop
(Chapter~3, Figure~3.4).

\section*{GA 4. Investigations, Experiments and Data Analysis}

Chapter~6 (Experimental Evaluation) is the primary evidence for this
outcome. Two designed experiments were run: a physical command-level test
across levels L0--L5 (41 tagged runs on the physical TurboPi), and a virtual
compositional-depth sweep (depths 4 to 120, three context-window settings,
270 runs in total). Every run was logged programmatically to CSV rather than
assessed by eye, with defined outcome and correctness metrics stated in
advance (Chapter~6, ``Evaluation Metrics''), and the results were analysed
quantitatively: per-depth and per-context breakdowns, and an explicit
comparison between the pipeline's own outcome label and the human
operator's recorded verdict where the two disagreed.

\section*{GA 5. Engineering Methods, Skills and Tools, Including Information
Technology}

The project required competent use of a wide tool stack: ROS2 (Humble)
running inside Docker on the TurboPi, a FastAPI backend, local LLM serving
via Ollama (Gemma~4~E4B as the primary planner and vision model, with
Qwen2.5-VL~3B evaluated as an alternative), a React/TypeScript frontend, and
a Python-based experiment harness for logging and analysing results.
Limitations of these tools were recognised and reported rather than hidden
-- for example, the open question of whether the Ollama server actually
enforced a reduced \texttt{num\_ctx} setting is stated explicitly as an
open item in Chapter~6, rather than silently assumed to be true.

\section*{GA 6. Professional and Technical Communication}

This outcome is demonstrated by the written report (this document) and by
the internal oral examination, both of which are required by the module
framework for this specific GA.

\section*{GA 8. Individual Work}

This project was completed through the work of an individual student,
with input from my advisor limited to our scheduled weekly meetings.
[[VERIFY -- state here, in your own words, your actual level of AI-tool
use in drafting the report, following the honest-disclosure precedent set
by other 2025 reports in this module.]]

\section*{GA 9. Independent Learning Ability}

The project required independent acquisition of knowledge well outside
undergraduate coursework: working with ROS2 topics, services, and
\texttt{rosbridge\_websocket} to drive the physical rover; understanding
local LLM serving behaviour and context-window/generation-length limits
well enough to design and interpret the depth sweep in Chapter~6; and
reverse-engineering the TurboPi's own ROS2 node graph and service interface
directly from its Docker image and source code, which was not documented
for this use case anywhere I could find. None of this was fully specified
in the original project brief and had to be learned independently over the
course of the project (Chapter~2 Literature Review; Chapter~5
Implementation).
